\documentclass[aps,prd,twocolumn,10pt,nofootinbib,superscriptaddress]{revtex4-2}
\usepackage{amsmath,amssymb,graphicx}
\usepackage{booktabs}
\usepackage[colorlinks=true,linkcolor=blue,citecolor=blue,urlcolor=blue]{hyperref}

\begin{document}

\title{Dark matter as a geometric scalar:\\
the Aether-Scalar-Tensor host for a cosmological-constant acceleration scale, and its price}
\author{C.\ Zimmerman}
\affiliation{Independent}
\date{June 2026}

\begin{abstract}
The hypothesis that galactic ``dark matter'' is a modification of dynamics with an acceleration scale fixed by the
cosmological constant, $a_0=c^2\sqrt{\Lambda/32\pi}=9.36\times10^{-11}\,\mathrm{m\,s^{-2}}$, requires a covariant host
that survives the gravitational-wave-speed bound, lenses correctly, and is consistent with the cosmic microwave
background. We argue that the Aether-Scalar-Tensor theory (AeST) is that host, and account honestly for what it buys
and what it costs. AeST is built from a single metric, a unit-timelike vector field (the ``aether''), and a
shift-symmetric scalar; it is ghost-free, propagates tensor modes at the speed of light in all backgrounds---surviving
GW170817 where the earlier bimetric TeVeS did not---and it yields two structural results we rely on: gravitational
lensing has no slip ($\Phi=\Psi$, so lensing mass equals dynamical mass) and the acceleration scale $a_0$ is
identically absent from linear cosmological perturbation theory. Crucially, in AeST the galactic scale $a_0$ and the
cosmological dark-matter-like component live in \emph{orthogonal} sectors of one free function: $a_0$ enters a forced
spatial-gradient term, while the cosmic ``dust'' that drives the third acoustic peak is a shift-symmetric k-essence
mode whose amplitude is a free integration constant. AeST therefore fits the full Planck spectrum---but at the price
of one early-universe density parameter that $a_0=\Lambda$ does not supply: the unification of the dark sectors is
galactic, not cosmological. We name the full price (a second density $\Omega$, a free cluster mass scale, an underived
$\mathcal{O}(1)$ coefficient $Z$) rather than hide it, and conclude that the geometric-scalar reading of dark matter
is viable and falsifiable, at one honest extra number.
\end{abstract}

\maketitle

\section{The thesis and the host}
If galaxy rotation curves reflect modified dynamics rather than unseen particles, and if the modifying acceleration
scale is set by the cosmological constant via $a_0=c^2\sqrt{\Lambda/32\pi}$~\cite{Milgrom1983,Skordis2021}, then the
proposal lives or dies on a relativistic completion. Three things have historically killed relativistic MOND: the
speed of gravitational waves (the bimetric Tensor-Vector-Scalar theory, TeVeS~\cite{Bekenstein2004}, predicted
$c_T\neq c$ and was excluded by GW170817~\cite{GW170817}); the under-prediction of gravitational lensing; and the
cosmic microwave background (CMB), whose third acoustic peak demands a clustering, pressureless component that a force
law does not provide. The Aether-Scalar-Tensor theory (AeST) of Skordis and Z\l o\'snik~\cite{Skordis2021} clears all
three. This paper takes AeST as the host for the cosmological-constant acceleration scale and accounts, both ways, for
what that costs.

\section{The geometric structure}
AeST is economical. Its field content is one metric $g_{\mu\nu}$, one unit-timelike vector $A_\mu$ with $A^\mu
A_\mu=-1$ (the ``aether,'' satisfying $\nabla_\mu A^\mu=3H$ on a Friedmann background), and one shift-symmetric scalar
$\phi$~\cite{Skordis2021}. Matter couples only to $g_{\mu\nu}$, so the weak equivalence principle holds exactly; the
vector $A_\mu$ selects a preferred cosmic rest frame, which is precisely the structure a de Sitter--Unruh reading of
inertia requires. The vector is the descendant of Einstein-Aether theory~\cite{Jacobson2001}; the disformal role it
plays in fixing lensing is inherited from TeVeS but realised here on a \emph{single} metric, with a kinetic term
engineered so that the tensor sector propagates at $c_T=c$ in every background. This is why AeST survived the
neutron-star merger that excluded TeVeS.

\section{Two structural results}
Two consequences of the $A_\mu$ structure are not tuned but forced, and we use both.

\paragraph{No-slip lensing.} Because the time component of the aether tracks the metric, $A^0\sim\sqrt{-g^{00}}$, the
two weak-field potentials are equal, $\Phi=\Psi$~\cite{Verwayen2024}. Light deflection ($\propto\Phi+\Psi$) and
dynamics ($\propto\Phi$) therefore read the same enhanced potential, and the lensing mass equals the dynamical mass
for any source. This repairs the historic under-lensing of non-relativistic MOND, and it means that cluster
strong$+$weak lensing, the Bullet cluster offset, and Einstein radii are not independent dark-matter detections but
restate the cluster dynamical residual.

\paragraph{The acceleration scale is CMB-invisible.} On a Friedmann background the spatial-gradient invariant of the
scalar vanishes identically, so the term that carries $a_0$ contributes nothing to the linear perturbation
equations~\cite{Skordis2021}. The galactic acceleration scale is therefore \emph{absent} from the CMB transfer
functions: it can neither help nor hurt the acoustic peaks. Whatever fits the CMB in AeST does so independently of
$a_0$.

\section{The decoupling: two orthogonal sectors}
The central structural fact, and the source of both the success and the cost, is that $a_0$ and the cosmological
dark-matter-like component occupy \emph{different} pieces of the AeST free function $\mathcal{F}(Y,Q)$. The galactic
scale enters only the spatial-gradient sector, through a forced fractional-power term $\mathcal{F}\supset
[2\lambda_s/3(1{+}\lambda_s)a_0]\,Y^{3/2}$ as $|\nabla\phi|\to0$---the $3/2$ power is the square-root force law, and it
is here, and only here, that $a_0$ appears~\cite{Skordis2021}. The cosmological component lives in the orthogonal
$\mathcal{K}(Q)$ sector: a shift-symmetric k-essence scalar with a minimum at $Q_0\neq0$ behaves, once integrated, as
pressureless dust, $\bar\rho\propto a^{-3}$, plus a cosmological constant. The dust amplitude is set by an integration
constant $I_0$, $8\pi\tilde{G}\bar\rho_0=Q_0 I_0$; in the authors' words the density ``is not (classically)
predicted.'' Galaxies and the early universe are thus controlled by independent slots of one function.

\section{The CMB at one extra number}
AeST reproduces the full Planck temperature and polarization spectra and the linear matter power spectrum to the
quality of $\Lambda$CDM---the first relativistic MOND theory to do so~\cite{Skordis2021}. The mechanism is the
$\mathcal{K}(Q)$ dust mode: a pressureless density $\propto a^{-3}$ that drives the acoustic peaks exactly as cold
dark matter does. We confirmed independently, with a standard Boltzmann code, that the third-peak height is a direct
measure of this clustering density: the ratio of third- to second-peak amplitude rises monotonically with the
physical cold density, from $\simeq0.53$ for a baryon-only universe to $\simeq0.98$ for the observed
$\Omega_c h^2\simeq0.12$, and a pure modified-inertia universe (in which $a_0$ adds nothing to the linear physics)
cannot reach the observed value by any tuning.

The price is explicit. The dust amplitude $I_0\sim\Omega_c h^2$ is a free integration constant, not delivered by
$a_0=c^2\sqrt{\Lambda/32\pi}$. The acceleration scale fixes the dark-\emph{energy} face ($\Omega_\Lambda=0.685$ for
the measured $\Lambda$), but the dark-energy and the CMB dust densities stand in ratio $\rho_{\rm DE}/\rho_c\simeq2.6$
today and redshift differently ($\,$const versus $(1{+}z)^3$, diverging by recombination), so one number cannot set
both. The slogan ``two dark sectors, one number'' holds for galaxy dynamics ($a_0\leftrightarrow\Lambda$ generates the
radial-acceleration and Tully--Fisher relations) but breaks at the CMB, which requires a second, independently tuned
density. The CMB is not evidence against the framework---AeST fits it---nor is it a flaw in the data; it is a cost in
unification economy of exactly one early-universe parameter.

\section{The cluster residual}
The same logic clarifies clusters. The dynamical residual at $R_{500}$, $\eta\simeq2.3$ on the eROSITA eRASS1
sample~\cite{Bulbul2024}, is shared with all modified dynamics and is not cured by $a_0=\Lambda$; by the no-slip
identity the cluster lensing residual is the \emph{same} number, not an independent one. Three honest points temper
it. First, the eRASS1 masses are weak-lensing-calibrated, and the lensing masses run $\sim110\%$ above the
hydrostatic and kinematic masses (which agree with each other), bracketing the true residual to $\eta\in[\sim1,2.3]$;
direct microcalorimeter measurements of non-thermal pressure and the projection bias of lensing masses pull the
convergent estimate to $\eta\sim1.3$--$1.9$. Second, the framework's own lower $a_0$ makes the residual $\sim13\%$
\emph{worse} ($\eta\propto a_0^{-1/2}$), so it is not a convention artifact. Third, the same $\mathcal{K}(Q)$ dust
that fits the CMB is the leading candidate to supply the residual at clusters, and the AeST scalar mass term yields a
right-signed enhancement above MOND at cluster scales~\cite{Durakovic2024}---but the mass term is non-monotone, is
squeezed between galaxy- and cluster-lensing constraints on the same parameter~\cite{Mistele2023}, and any such cure
introduces its own density scale, not $a_0=\Lambda$. The cluster residual is real, soft, systematic-limited, and
uncured as a unification win.

\section{The price, named}
We state the full cost of the geometric-scalar reading, rather than bury it. Beyond $a_0=c^2\sqrt{\Lambda/32\pi}$ the
AeST host requires: (i) one early-universe density, the integration constant $I_0\sim\Omega_c h^2$, that drives the
CMB and is the leading cluster candidate---orthogonal to $a_0$; (ii) a free cluster-scale mass $\mu$ in the scalar
sector, constrained from opposite directions by galaxy and cluster lensing; (iii) the shape of the free function
$\mathcal{F}(Y,Q)$, chosen to fit; and (iv) the dimensionless coefficient $Z=\sqrt{32\pi/3}\simeq5.79$ in
$a_0=cH_\Lambda/Z$, which the de Sitter geometry fixes only to $\mathcal{O}(1)$ and which we do not assert as derived.
What the geometry \emph{does} fix, with no freedom, is the form and scale: $a_0\sim c^2\sqrt{\Lambda}$ is the
Gibbons--Hawking horizon acceleration, and the appearance of $a_0$ in a $Y^{3/2}$ term reproduces the square-root
force law. Form and scale: forced. Coefficient and cosmological unification: posited.

\section{Conclusion}
The proposition that dark matter is an illusion---a geometric scalar field in a Lorentz-violating completion of
general relativity, with an acceleration scale set by the cosmological constant---has, in AeST, a host that is
ghost-free, gravitational-wave-safe, lenses without slip, and fits the cosmic microwave background. That is a
substantive position, not a slogan. It is also not free: the cosmological dark-matter density is one parameter
distinct from the galactic $a_0=\Lambda$ unification, the cluster residual is real if soft, and the coefficient $Z$
is an honest posit. The geometric-scalar reading should be judged on those terms---one extra early-universe number,
and a falsifiable, distinctive prediction the particle picture cannot make: a redshift-declining
$a_0(z)\propto\sqrt{\rho_{\rm DE}(z)}$.

\begin{acknowledgments}
The Boltzmann third-peak recomputation, the eRASS1 residual at the framework's own $a_0$, and the AeST sector
accounting were performed and cross-checked numerically; the AeST identities are quoted from the primary literature.
Both-ways audit ledgers are archived with this work.
\end{acknowledgments}

\begin{thebibliography}{99}
\bibitem{Milgrom1983} M.~Milgrom, Astrophys.\ J.\ \textbf{270}, 365 (1983).
\bibitem{Skordis2021} C.~Skordis, T.~Z\l o\'snik, Phys.\ Rev.\ Lett.\ \textbf{127}, 161302 (2021) [AeST].
\bibitem{Bekenstein2004} J.~D.~Bekenstein, Phys.\ Rev.\ D \textbf{70}, 083509 (2004) [TeVeS].
\bibitem{GW170817} B.~P.~Abbott \emph{et al.} (LIGO/Virgo, Fermi, INTEGRAL), Astrophys.\ J.\ \textbf{848}, L13 (2017).
\bibitem{Jacobson2001} T.~Jacobson, D.~Mattingly, Phys.\ Rev.\ D \textbf{64}, 024028 (2001) [Einstein-Aether].
\bibitem{Verwayen2024} P.~Verwayen, C.~Skordis, C.~Z\l o\'snik, Mon.\ Not.\ R.\ Astron.\ Soc.\ \textbf{531}, 272 (2024).
\bibitem{Bulbul2024} E.~Bulbul \emph{et al.}, Astron.\ Astrophys.\ \textbf{685}, A106 (2024) [eRASS1].
\bibitem{Durakovic2024} A.~Durakovi\'c, C.~Skordis (2024) [AeST scalar mass term in clusters].
\bibitem{Mistele2023} T.~Mistele, S.~McGaugh, S.~Hossenfelder, Astron.\ Astrophys.\ \textbf{676}, A100 (2023).
\end{thebibliography}

\end{document}
